\section{第~\ref{sec:acet}~章：加入\nt{ACET}}

乙酰胆碱的三种作用已在脑片和活体脑中测得，写入与读出的冲突在本书模型中得到展示，而\nt{ACET}充当写使能线并报告预期的不确定性，是有部分实验支持的假说。

{\footnotesize
\setlength{\tabcolsep}{3pt}\renewcommand{\arraystretch}{1.15}
\begin{longtable}{>{\raggedright}p{0.5cm}>{\raggedright}p{4.0cm}>{\raggedright}p{5.6cm}>{\raggedright}p{2.3cm}>{\raggedright\arraybackslash}p{1.7cm}}
\toprule
序 & 命题 & 实验与测量内容 & 来源 & 状态 \\
\midrule\endfirsthead
\toprule
序 & 命题 & 实验与测量内容 & 来源 & 状态 \\
\midrule\endhead
1 & 脑中的\nt{ACET}神经元不多，集中在几个核团中，而输出遍布整个皮层：这是广播通道 & 大鼠，乙酰胆碱合成酶抗体染色和逆行标记：皮层、海马、杏仁核、嗅球和丘脑的乙酰胆碱主要来自基底前脑和脑干的六个紧凑细胞群 & \cite{Mesulam1983} & 已测量 \\
2 & 皮层中乙酰胆碱水平在清醒时高，在慢波睡眠中低 & 自由活动猫的皮层和海马微透析并同步记录脑电：安静清醒时乙酰胆碱释放比慢波睡眠高$\sim100\%$，活跃清醒时高$\sim175\%$；快速眼动睡眠中皮层的释放与清醒时相当 & \cite{Marrosu1995,Hasselmo1999} & 已测量 \\
3 & 信号的含义由受体决定：乙酰胆碱既有快速的通道型受体，也有在细胞内启动反应的慢速受体（命题~\ref{prop:receptor}） & 测量综述：乙酰胆碱对兴奋性、传递和可塑性的作用取决于释放部位、受体亚型（烟碱型为离子通道，毒蕈碱型经 G 蛋白）和接收细胞 & \cite{Picciotto2012} & 已测量 \\
4 & 第一种作用：削弱内部连接，几乎不影响外部输入（\eqref{eq:acetnet}中的因子$1-a$） & 大鼠嗅皮层脑片：$100$\,\textmu M 乙酰胆碱激动剂使同一细胞中内部连接（Ib 层）的电位下降$60\%$以上，外部输入（Ia 层）的下降不到$15\%$；抑制发生在突触前。海马脑片（CA1）：内部通路和输入通路分别为$89\%$和$40\%$ & \cite{HasselmoBower1992,HasselmoSchnell1994} & 已测量 \\
5 & 第二种作用：增强输入（因子$\frac{1+a}{2}$） & 新皮层脑片：烟碱型受体增强来自丘脑的输入突触，对皮层内突触无作用；毒蕈碱型受体削弱两类突触。乙酰胆碱提高神经元兴奋性（综述） & \cite{Gil1997,Hasselmo2006} & 已测量 \\
6 & 第三种作用：使权重更容易改变（速率$\eta a$） & 大鼠：电刺激基底核（皮层乙酰胆碱的来源）并与纯音配对，使成年动物的听觉皮层图谱按所呈现的声音发生显著重组 & \cite{KilgardMerzenich1998,Hasselmo2006} & 已测量 \\
7 & \eqref{eq:acetnet}第二个方程中针对稀疏模式的赫布规则给出可根据部分补全模式的联想记忆 & 采用稀疏模式和协方差规则的网络容量计算：活跃神经元比例$p$小时，网络存储的模式明显多于$p=1/2$时 & \cite{Tsodyks1988} & 理论 \\
8 & 低$a$时的写入损坏记忆：网络写下的是回忆出的而不是看到的；$a=0.1$时的损坏不比$a=0$时轻 & \texttt{models/\allowbreak acet\_memory.py}：$200$个神经元，五个各含$20$个神经元的模式，新模式与其中一个共有$10$个神经元，$10$次运行。$a=0$时新模式未被记住，旧模式带出$10$个外来神经元中的$5.9$个；$a=0.1$时新模式可回忆（$0.97$），但两者融合：新模式带出$7.3$个外来神经元，旧模式带出$7.9$个；从$a=0.2$起不到一个 & 本章推导 & 模型 \\
9 & 命题~\ref{prop:writeread}：不存在使写入和读出同样良好的$a$值（图~\ref{fig:acetsim}） & 同一脚本，学习速率$\eta a$：$a\ge0.9$时新模式一次呈现即被完整记住，$a=0.75$时为$0.8$，$a\le0.5$时记不住。凭一半模式读出：$a\le0.2$时为$1.0$，$a=0.3$时为$0.96$，$a=0.4$时为$0.04$ & 本章推导 & 模型 \\
10 & 在脑中，一根广播导线——\nt{ACET}——切换写入和读出 & 这一想法来自根据脑片测量建立的模型。最直接的实验在人类身上：毒蕈碱受体阻断剂东莨菪碱损害新词对的学习，对与旧词对重叠最多的影响最大，但不损害已学词对的回忆。没有对网络两种模式的直接测量 & \cite{HasselmoAndersonBower1992,Atri2004} & 假说 \\
11 & 滤波器~\eqref{eq:kalman}是最优的；输入权重和学习速率是同一个量$P/R$，稳态下$P=\sqrt{qR}$，记忆为$\sqrt{R/q}$（命题~\ref{prop:kalman}） & 不需要实验：卡尔曼—布西滤波器的最优性是估计理论的经典结果；稳态在本章推导 & \cite{KalmanBucy1961} & 理论 \\
12 & \nt{ACET}向网络报告一个类似$P$的量——预期的不确定性 & 在一个注意的概率模型中提出。没有乙酰胆碱水平跟随估计不确定性的直接测量 & \cite{YuDayan2005} & 假说 \\
13 & 部分\nt{ACET}神经元在二十毫秒左右内对奖励和惩罚作出响应，强化越出乎意料，响应越强 & 小鼠，经光遗传鉴定的基底前脑胆碱能神经元：对奖励和惩罚的响应潜伏期为$18\pm3$\,ms，幅度随强化的意外程度增大，两个核团的神经元响应相似 & \cite{Hangya2015} & 已测量 \\
14 & \nt{NORA}发现世界的意外变化，\nt{ACET}使网络保持在写入模式 & 这种分工是在模型中提出的。间接证据：人类与\nt{NORA}相关的瞳孔直径跟踪变化和学习速率。没有对\nt{NORA}--\nt{ACET}这一对的直接检验 & \cite{YuDayan2005,Nassar2012} & 假说 \\
15 & 在慢波睡眠中，处于读出模式的网络回放白天记录的内容，这些回放把它转写到长期记忆中 & 大鼠海马神经元集群记录：白天一起工作的细胞在随后的慢波睡眠中更常一起发放，且顺序相同——已测量。关于转写：学习后抑制海马涟漪会损害空间记忆，延长自发涟漪则改善它；原因是\nt{ACET}水平低，这一点未经证明 & \cite{WilsonMcNaughton1994,Skaggs1996,Girardeau2009,FernandezRuiz2019,Hasselmo1999} & 假说 \\
\bottomrule
\end{longtable}}
